% !TeX root = ../main.tex
\section{Problem Formulation and Benchmark Design}
\label{sec:method}

\subsection{Problem Formulation}
The road network is a directed graph $G=(V,E)$, where $V$ is the set of intersections and $E$ is the set of road segments. Each segment $e_i \in E$ has a feature vector $\mathbf{x}_i \in \mathbb{R}^{d}$, a length $\ell_i$, a free-flow travel time $t_i$, and a flood label $y_i \in \{0,1\}$. The benchmark has two tasks.

\subsubsection{Classification task}
A model estimates the probability that segment $e_i$ floods as follows (Eq.~\ref{eq:prob}).
\begin{equation}
  p_i = P(y_i = 1 \mid \mathbf{x}_i)
  \label{eq:prob}
\end{equation}
For the graph model, the input also includes the features of neighboring segments, taken from the line graph $L(G)$, whose nodes are the segments of $G$ and whose edges join segments that share an intersection.

\subsubsection{Routing task}
The probabilities are converted into edge costs as below (Eq.~\ref{eq:cost}).
\begin{equation}
  C_i = t_i \,(1 + \lambda\, p_i)
  \label{eq:cost}
\end{equation}
Here, $\lambda \geq 0$ controls risk aversion. For an origin--destination (OD) pair $(o,d)$ with the set of paths $\Pi_{od}$, the flood-aware route is given below (Eq.~\ref{eq:route}).
\begin{equation}
  \pi^{*} = \arg\min_{\pi \in \Pi_{od}} \sum_{e_i \in \pi} C_i
  \label{eq:route}
\end{equation}
A route $\pi$ is evaluated by its observed flood exposure and its travel-time increase over the shortest-time route $\pi_0$ as below (Eq.~\ref{eq:routemetrics}).
\begin{equation}
  F(\pi) = \sum_{e_i \in \pi} y_i\, \ell_i, \qquad
  \Delta T(\pi) = \frac{T(\pi) - T(\pi_0)}{T(\pi_0)}
  \label{eq:routemetrics}
\end{equation}
Here, $T(\pi)=\sum_{e_i \in \pi} t_i$. The research question is whether the model with the best classification score also gives the lowest $F$ at a given $\Delta T$.

\subsection{Models for Comparison}
Six models from five families are compared, each with fixed hyperparameters:
\begin{itemize}
  \item \textbf{Linear:} logistic regression (LR) with the liblinear solver, a transparent baseline.
  \item \textbf{Kernel:} support vector machine (SVM) with a radial basis function kernel~\cite{cortes1995}, $C=1$, and Platt scaling.
  \item \textbf{Tree ensembles:} random forest (RF)~\cite{breiman2001} with 100 trees, and XGBoost~\cite{chen2016} with 100 trees of depth 5 and a learning rate of 0.08.
  \item \textbf{Neural:} multilayer perceptron (MLP) with two hidden layers (64 and 32 units), ReLU activations, and dropout (0.2), trained full-batch with Adam (learning rate 0.01, weight decay $10^{-4}$) for 80 epochs.
  \item \textbf{Graph:} GraphSAGE~\cite{hamilton2017} on the line graph $L(G)$, with two SAGEConv layers (32 hidden channels), mean neighborhood aggregation, and the same Adam settings trained for 100 epochs.
\end{itemize}
All models receive the same features and folds. Class weights handle class imbalance (balanced weights for LR, the SVM, and RF, and a positive-class weight for XGBoost, the MLP, and GraphSAGE). Features are standardized within each training fold for all models except the tree ensembles; GraphSAGE, which is trained transductively, standardizes over all segments. Hyperparameters are not tuned.

\subsection{Reproducibility}
\label{sec:repro}
All random operations use the fixed seed 42: the fold assignment, the model initialization, and the sampling of OD pairs. Spatial folds are built by shuffling the 500~m blocks (Section~\ref{sec:cv}) with this seed and assigning them to the five folds in turn, and random folds use a shuffled five-fold split with the same seed. Experiments ran on \TODO{CPU model, RAM} under \TODO{operating system} with \TODO{Python version}, PyTorch 2.14, XGBoost 3.2, and PyTorch Geometric 2.8. The road graph, the feature table with labels, the feature-extraction scripts, and the notebook that reproduces the reported tables and figures are available in an anonymous repository: \TODO{anonymous repository link}. The raw simulation output, DEM, and rain-gauge records are not included, but the feature table holds every model input and label.

\subsection{Classification Benchmark}
\label{sec:cv}
Neighboring segments share terrain and drainage, so random splits leak information between training and test sets~\cite{roberts2017}. The district is therefore divided into square blocks of 500~m, which are assigned to five folds and held out in turn. For GraphSAGE, held-out segments stay in the graph, but their labels are hidden during training. Random five-fold cross-validation is also run. The optimism gap of model $m$ is defined as below (Eq.~\ref{eq:gap}).
\begin{equation}
  \Delta_m = \text{AUC}^{\text{random}}_m - \text{AUC}^{\text{spatial}}_m
  \label{eq:gap}
\end{equation}
The primary metric is the area under the precision--recall curve (AUC-PR), whose no-skill baseline equals the share of flooded segments (0.425). AUC-ROC, F1, and recall on the flooded class are also reported, and calibration is measured with the Brier score~\cite{brier1950}, because the probabilities become routing costs. F1 and recall use a threshold of 0.5. Differences in AUC-PR across folds are tested with the Friedman test and pairwise Wilcoxon signed-rank tests against the best model~\cite{demsar2006}; with five folds, the smallest attainable two-sided Wilcoxon $p$-value is 0.0625, so these tests show trends only. Feature effects are explained with TreeSHAP~\cite{lundberg2017} on an XGBoost model fitted to all segments. To compare models, each feature is ranked by a model-specific importance: the absolute coefficient for LR, the impurity importance for RF, the mean absolute SHAP value for XGBoost, and, for the SVM, MLP, and GraphSAGE, the absolute correlation between the feature and the model's out-of-fold probabilities.

\subsection{Routing Benchmark}
Every segment receives the probability predicted in the fold where it was held out, so no model routes on labels it was trained on. Free-flow travel times use speeds by road class, from 50~km/h on primary roads to 20~km/h on residential streets. For 150 random OD pairs at least 300~m apart, routes are computed with Dijkstra's algorithm for each model and for two baselines: shortest time ($p_i=0$) and an oracle ($p_i=y_i$) that knows the simulated labels. The main results use $\lambda=5$; because the models output probabilities on different scales, the same $\lambda$ makes them detour by different amounts, so $\lambda$ is also varied over $\{1, 2, 5, 10, 20\}$ on the first 50 OD pairs to compare models at equal travel-time cost. Finally, the model rankings by classification and by routing are compared.

\subsection{Ablation, Sensitivity, and Efficiency}
The feature-group ablation removes the terrain, hydrological, rainfall, or network features in turn for XGBoost and GraphSAGE. For GraphSAGE, the number of layers is varied, and a version without graph edges isolates the effect of connectivity. Sensitivity is tested with XGBoost for the block size (250, 500, and 1000~m) and for the labeling of flooded-but-passable segments. Efficiency is measured by the number of parameters, the training time, the inference time for the district, and the routing time per OD pair, on the hardware described in Section~\ref{sec:repro}.
